\documentclass{syllabus}

\begin{document}

\thispagestyle{fancy}
\date{Spring 2026}
\author{}
\title{PHY (EGR) 3120: {Electronics}}
\maketitle

\meeting{Lecture: & MWF, 1:35-2:40 PM  \\ &Lab period: & T, 12:30-2:10 PM}
\officehours{Tuesday: & 2:10-3:00\\ &Wednesday: & 2:40-3:30}
\textbook{Basic Electronics for Sci. and Eng. (2/e)}{Eggleston}{Cambridge}{2025}
\prerequisites{PHY 2110: & (Physics II) or\\ & PHY 2210: & (General Physics II)}

\section*{General Information}

\infotable
\hspace{\fill}
\includegraphics[height=2.25in]{Assets/electronicsCover.png}
\hspace{.25in}





\section*{Course overview}

Introduction to analog and digital circuits. Topics include passive and active analog electronic components, DC and AC circuit analysis, amplifiers, filters, binary and digital systems, logic gates, and microcontroller programming. 

Electronics form the basis of much of modern technology. This course will give an introduction to the fundamental principles that are used in most electronics applications ranging from power electronics, consumer electronics, computer electronics, etc. 

Topics in this course include
\begin{enumerate}[label=\textcolor{CarthageDark}{\bfseries\arabic*.}]
    \item Prototyping circuits using a breadboard.
\item Using common electronics test equipment (multimeter, signal generator, oscilloscope, etc.) in the circuit prototyping process.
\item Reading, interpreting, and drawing electronics schematics using appropriate symbols and layout conventions
\item A thorough theoretical and practical exploration of passive circuits (circuits made up of linear elements such as capacitors, resistors, and inductors) in both the time and frequency domain.
\item An introduction to semiconductor theory and its application to basic semiconductor devices  such as diodes, bipolar junction transistors (BJTs), and field effect transistors (FETs).
\item An introduction to digital logic, binary arithmetic, and sequential logic.
\item Designing circuits which utilize microcontrollers and programming microcontrollers to respond appropriately.

\end{enumerate}


\subsection*{Student learning outcomes} 

\noindent
Upon successful completion of this course, the student will be able to:
\begin{enumerate}[label=\textcolor{CarthageDark}{\bfseries\arabic*.}]
    \item Prototype analog circuits using breadboards, power supplies, and function generators and probe the behavior of these circuits using multimeters and oscilloscopes.
\item Explain how capacitors, resistors, inductors, diodes, and transistors function in an electronic circuit as well as the physical processes underlying these functions.
\item Apply op-amps to a variety of analog use cases.
\item Perform basic logic and binary arithmetic operations and design basic sequential digital systems.
\item Apply understanding of microcontrollers to data collection and control tasks.
\end{enumerate}

\textbf{Patience} is defined as the quality by which one endures of prolonged, arduous, and difficult things for the sake of some good. This course will provide opportunities to practice and reflect on the virtue of patience. Cf. the engineering department value of \textbf{tenacity}.

\section*{Course components}

\paragraph*{Laboratory and in-class exercises: }
The laboratory activities form the backbone of this course. The standard format for laboratory activities will be to build and characterize the behavior of a series of circuits. See the laboratory notebook section below for more instructions on the laboratory deliverables. These activities are an essential complement to the lecture and theoretical portion of this course, as you will be analyzing the working principles of various common circuits in light of theoretical considerations learned in lecture. Laboratory exercises are expected to be accomplished during the lab period: if you need additional time outside of class, this will have to be arranged with the instructor.

In addition to the laboratory exercises, there will be occasional \textbf{in-class exercises}. These will be shorter in format and largely graded on a pass-fail basis.

\paragraph{Homework} 
Homework exercises will be assigned sporadically throughout the course as new topics arise in class. These are designed to give a more detailed theoretical analysis of some circuit and device properties without the aid of a physical instantiation present to test. 

\paragraph{Projects}
\begin{enumerate}[label=\textcolor{CarthageDark}{\bfseries\arabic*.}]
\item \textbf{Analog project:} Design and build a circuit that receives analog input from the outside world (or from an oscillator) and transforms the signal into a more useful output.
\item \textbf{Digital project:} design a finite state machine and implement it in circuitverse. 
\item \textbf{Microcontroller project:} Design and build a circuit that receives some input from the outside world and uses this information to control digital components. This will also involve programming a microcontroller to control this process.
\end{enumerate}

\paragraph{Exams} 
Two mid-term exams (on passive and active analog circuits) and a final exam (on digital and sequential logic) will be administered during the semester. These will be 65 minutes in length and may be accompanied by a lab practical section.

\paragraph{Reflection} 
Weekly reflection exercises will be completed following each lab to encourage engagement with class activities as an opportunity to grow in the virtue of patience as outlined above.


\section*{Grading Scheme}
\subsection*{Component Weights}

\noindent\begin{tabular}{rr}
\textbf{Component}      & \textbf{Weight}       \\ \hline
\textbf{Homework}       & 12\%  \\ \\
In-class exercises       & 6\%       \\
\underline{+ Laboratory exercises}       & \underline{+ 12\%}       \\
\textbf{Lab total}       & 18\%       \\ \\
Analog project       & 12\%       \\
Digital project       & 10\%       \\
\underline{+ Microcontroller project}       & \underline{+ 12\%}       \\
\textbf{Project total}       & 34\%       \\ \\
\textbf{Exams}          & 3 $\times$ 10\% \\
                        & (30\%) \\ \\
\textbf{Reflection}     & 6\%
\end{tabular}%



\lettergrades{a) number of extra credit assignments adequately completed, b) a pattern of seeking help in the course at office hours, c) number of missing homework assignments.}

\section*{Tips for success}
\textbf{Before lab}: Look over the circuits before lab and try to wrap your head around how they work/what they do. Complete the prelab activities. 
\textbf{Participate fully in the lab} and other in-class activities: make sure you understand what is going on in the laboratory exercises. Don’t let your partner do all the heavy lifting. For example, make sure you take turns using the oscilloscope.


%\tipsforsuccess{upper}

\section*{Policies}
\subsection{General course policies}
Expect a response within one business day for emails sent to the instructor. To ensure a prompt response, format the subject line as: [PHY 3120] *insert subject here*

All required material will be covered in lecture, but can also be found in the parallel readings, so in the case of absence students should consult the scheduled readings. Students found sleeping will be brought to the attention of the class for public ridicule and no effort made to wake them up.

Students are expected to arrive and leave on time. I will attempt to arrive five minutes early and be available for at least five minutes after class.

\latework{Easy}
\collegepolicies

\end{document}
